% Part: incompleteness
% Chapter: theories-computability
% Section: intepretability

\documentclass[../../../include/open-logic-section]{subfiles}

\begin{document}

\olfileid[hi]{inc}{tcp}{itp}

\olsection{जिन सिद्धांतों में $\Th{Q}$ व्याख्येय है वे अनिर्णेय हैं}

हम इन परिणामों को और भी प्रबल बना सकते हैं। अनौपचारिक रूप से, किसी भाषा
$\Lang{L_1}$ की किसी दूसरी भाषा $\Lang{L_2}$ में व्याख्या करने में
$\Lang{L_2}$ के !!{formula}s द्वारा $\Lang{L_1}$ का विश्व, संबंध-चिह्न और
फलन-चिह्न परिभाषित करना शामिल है। यद्यपि हम इसके लिए समय नहीं देंगे, इस
परिभाषा को सटीक बनाया जा सकता है।

\begin{thm}
  मान लें $\Th{T}$ ऐसी भाषा का सिद्धांत है जिसमें अंकगणित की भाषा की ऐसी
  व्याख्या की जा सकती है कि $\Th{T}$,~$\Th{Q}$ की व्याख्या के साथ संगत हो।
  तब $\Th{T}$ अनिर्णेय है। यदि $\Th{T}$,~$\Th{Q}$ के अभिगृहीतों की
  व्याख्या को सिद्ध करता है, तो $\Th{T}$ का कोई संगत विस्तार निर्णेय नहीं
  है।
\end{thm}

प्रमाण पिछले प्रमेय के प्रमाण का केवल एक छोटा संशोधन है; प्रतिवाद का उपयोग
करके $\Th{Q}$ और $\Th{\bar Q}$ का वियोजन प्राप्त किया जा सकता है।
चयन-अभिगृहीत सहित ज़र्मेलो--फ्रेंकेल समुच्चय-सिद्धांत $\Th{ZFC}$ को ऐसा
अभिगृहीती आधार माना जा सकता है जो सामान्य गणित का बड़ा भाग करने के लिए
पर्याप्त शक्तिशाली है। $\Th{ZFC}$ में प्राकृतिक संख्याएँ परिभाषित की जा
सकती हैं, और इस व्याख्या के माध्यम से $\Th{Q}$ के अभिगृहीत सत्य हैं। अतः
हमारे पास है:

\begin{cor}
$\Th{ZFC}$ का कोई निर्णेय विस्तार नहीं है।
\end{cor}

\begin{cor}
$\Th{ZFC}$ का कोई पूर्ण, संगत, संगणनीयतः !!{axiomatizable} विस्तार नहीं है।
\end{cor}

$\Th{ZFC}$ की भाषा में केवल एक द्वि-स्थानीय संबंध $\in$ है। (वास्तव में
समता की भी आवश्यकता नहीं है।) अतः हमारे पास है:

\begin{cor}
किसी द्वि-स्थानीय संबंध-चिह्न वाली प्रत्येक भाषा का प्रथम-कोटि तर्कशास्त्र
अनिर्णेय है।
\end{cor}

यह परिणाम दो एक-स्थानीय फलन-चिह्नों वाली प्रत्येक भाषा तक विस्तृत होता है,
क्योंकि उनका उपयोग करके द्वि-स्थानीय संबंध-चिह्न का अनुकरण किया जा सकता है।
अभी उद्धृत परिणाम सटीक सीमाएँ देते हैं: पता चलता है कि केवल
\emph{एक-स्थानीय} संबंध-चिह्नों और अधिकतम एक \emph{एक-स्थानीय} फलन-चिह्न
वाली भाषा के लिए प्रथम-कोटि तर्कशास्त्र निर्णेय है।

एक और रोचक तथ्य। हम जानते हैं कि मानक प्रतिरूप में सत्य, भाषा
$\Obj 0$, $'$, $+$, $\times$, $<$ के वाक्यों का समुच्चय अनिर्णेय है।
वास्तव में अन्य चिह्नों के पदों में $<$ परिभाषित किया जा सकता है, और फिर
$\times$ तथा $'$ के पदों में $+$ परिभाषित किया जा सकता है। अतः भाषा
$\Obj 0$, $'$, $\times$ में सत्य वाक्यों का समुच्चय अनिर्णेय है। दूसरी ओर
प्रेस्बर्गर (Presburger) ने दिखाया कि अंकगणित की भाषा में सत्य, भाषा
$\Obj 0$, $'$, $+$ के वाक्यों का समुच्चय निर्णेय है। किंतु वह प्रक्रिया
संगणनात्मक रूप से अव्यावहारिक है।

\end{document}
